\chapter{Reference Polynomial Definitions}
\label{app:polynomial-definitions}

This appendix collects the lower-level polynomial definitions used as reference examples throughout the thesis. They are not built into the compiler. Rather, they are small benchmarks for the expressiveness of the DSL: the same language constructs can be combined to describe other polynomial expressions.

Let \(M=(M_1,\ldots,M_B)\) be a sequence of field elements in \(\mathbb{F}_p\), where \(p=2^\pi-\theta\).

\paragraph{MMH.}
For keys \(r_1,\ldots,r_{B-1}\), the MMH-style lower-level hash is
\[
    \operatorname{MMH}(r,M)
      = \sum_{i=1}^{B-1} M_i r_i + M_B.
\]
It is represented in the DSL by a sum of products plus the final message element.

\paragraph{SQH.}
For keys \(r_1,\ldots,r_{B-1}\), the SQH-style hash replaces multiplication by squaring:
\[
    \operatorname{SQH}(r,M)
      = \sum_{i=1}^{B-1} (M_i+r_i)^2 + M_B.
\]
This is the motivating example for a dedicated square operation in the IR and backend.

\paragraph{NMH.}
For keys \(r_1,\ldots,r_B\), the NMH-style hash groups elements pairwise:
\[
    \operatorname{NMH}(r,M)
      =
      \sum_{i=1}^{\lfloor B/2 \rfloor}
        (M_{2i-1}+r_{2i-1})(M_{2i}+r_{2i})
      + (B \bmod 2)M_B.
\]
The final term is present only when \(B\) is odd.

\paragraph{HKM.}
The HKM definition is recursive. The paper version appears to contain a typo in the recursive call; the intended recurrence is the one shown here. For \(B=1\),
\[
    \operatorname{HKM}((r_1),(M_1)) = r_1 + M_1.
\]
For larger \(B\), let \(h=\lfloor B/2\rfloor\). Then
\[
\begin{aligned}
    \operatorname{HKM}((r_i)_{i=1}^{h+1},(M_i)_{i=1}^{B})
    &=
    \operatorname{HKM}((r_i)_{i=1}^{h},(M_i)_{i=1}^{2h-1})  \\
    &\quad{}\cdot (r_{h+1}+M_{2h}) + (B \bmod 2)M_B.
\end{aligned}
\]
For generated code we use an equivalent iterative formulation expressed as a left fold, because this avoids recursive generated C and exposes carry placement to the compiler.
Using zero-based indices, the iterative version used by the DSL module is:
\[
\begin{aligned}
    A_0 &= r_0 + M_0, \\
    A_{j+1}
        &= A_j \cdot (r_{j+1}+M_{2j+1}) + M_{2j+2},
        \qquad 0 \leq j < \left\lfloor\frac{B-1}{2}\right\rfloor .
\end{aligned}
\]
For odd \(B\), the final value is \(A_{\lfloor(B-1)/2\rfloor}\). For even \(B\), one additional multiplication is required:
\[
    A_{\lfloor(B-1)/2\rfloor}\cdot (r_{B/2}+M_{B-1}).
\]
This is exactly the shape expressed by the \texttt{foldl} implementation: each loop iteration consumes two message elements and one new key element, and the even-length case is handled after the fold.

\chapter{IR Instruction Reference}
\label{app:ir-instructions}

This appendix summarizes the IR instructions used by the compiler. Instructions operate on IR operands and produce or update an IR temporary. The IR type of the result determines whether an instruction is scalar, a single field element, or a vectorized collection of field elements.

\begin{description}
    \item[\texttt{const}] Creates a constant operand. Used for integer constants and zero-initialized field elements.
    \item[\texttt{copy}] Copies one operand into an existing temporary. Used primarily for accumulator updates and branch results.
    \item[\texttt{call}] Calls a helper function and stores its return value.
    \item[Index arithmetic] \texttt{add}, \texttt{sub}, \texttt{mul}, \texttt{div}, and \texttt{mod} compute scalar arithmetic on index operands. For field operands, \texttt{add} and \texttt{mul} denote field addition and field multiplication.
    \item[Comparisons] \texttt{eq}, \texttt{neq}, \texttt{lt}, \texttt{le}, \texttt{gt}, and \texttt{ge} compare scalar values. They return index values used as conditions.
    \item[\texttt{select}] Constant-time selection between two already computed operands, controlled by a scalar condition.
    \item[\texttt{load}] Loads one field element from a byte buffer and decodes it into scalar limb representation.
    \item[\texttt{vload}] Loads several consecutive field elements from a byte buffer and decodes them into vectorized limb representation.
    \item[\texttt{square}] Computes a scalar field square using the dedicated squaring algorithm.
    \item[\texttt{carry}] Performs scalar carry propagation and partial Crandall reduction on one field element.
    \item[\texttt{full\_reduction}] Performs the final scalar reduction that ensures the stored field element is below the field prime.
    \item[\texttt{vadd}] Adds vectorized field elements lane-wise.
    \item[\texttt{vmul}] Multiplies vectorized field elements lane-wise.
    \item[\texttt{vsquare}] Squares vectorized field elements lane-wise.
    \item[\texttt{vcarry}] Performs carry propagation lane-wise on vectorized field elements.
    \item[\texttt{splat}] Replicates a scalar field element into every SIMD lane.
    \item[\texttt{lane\_powers}] Builds a vectorized field element whose lanes contain specified powers or multiples of a scalar field element. This is used by vectorized Horner lowering.
    \item[\texttt{horiz\_add}] Horizontally adds the lanes of a vectorized field element, producing one scalar field element.
\end{description}
